\par\noindent\SourceLeaf{55}%
\hspace*{\parindent}So scheint sich ein sehr verwickeltes Gesetz mittels der Binomialkoeffizienten sehr kurz ausdrücken zu lassen.
Denn wenn man diese Koeffizienten, wie es möglich ist, in Produkte entwickelt, kürzen sich zwar einige Faktoren des einen gegen Faktoren des anderen, doch wird der so entstehende Ausdruck viel länger als der frühere. Um das am Ergebnis zu prüfen, geben wir für den beliebigen $r^{\text{ti}}$ Koeffizienten das allgemeine Glied $h^{\text{tus}}$ auf diese Weise in möglichst kurzer Form an

\[
\frac{\nu.(\nu-1).\,.\,.(\nu-(r+2h-1))}{1.2^2.3^2.\,.\,. h^2} . \frac{1}{(h+1)(h+2).\,.\,.(h+r)}\frac{1}{2^{r+2h-1}}n^{r+2h}.
\]

\textls[120]{Anmerkung.}
Wer das allgemeine Gesetz aus den Reihen ableiten will, die \textsc{Euler}\HistoricalNote{*)}{*)}{\textsc{Euler}, a. a. O., Problem 36.} für einige Anfangskoeffizienten angegeben hat, wird eine solche Formel erhalten. Der Weg, den der berühmte Verfasser zur Untersuchung dieses Gesetzes angibt, scheint für die Gewinnung der bereits dargelegten Reihen zwar weniger direkt und weniger einfach, erscheint aber in anderer Hinsicht ausgezeichnet und keineswegs zu vernachlässigen.
Das ganze Verfahren besteht darin, für die Größe $(1+n\cos\varphi)^{\nu}$ die Reihe

\[
A+A^{(1)}\cos\varphi+A^{(2)}\cos2\varphi+.\,.\,.\,.
\]

\noindent anzusetzen und durch Differentiation das Rekursionsschema herzuleiten, das die Abhängigkeit ihrer Koeffizienten voneinander ausdrückt. Da es nur drei Glieder umfasst, liefert es jeden beliebigen Koeffizienten aus den beiden jeweils vorhergehenden. Liegen also die ersten beiden in endlicher Form vor, so lassen sich auch alle übrigen in endlicher Form angeben. Diese Beziehung zwischen den Koeffizienten lässt sich zwar aus den bereits gefundenen Reihen herleiten, doch ergibt sie sich leichter mit \textsc{Euler}s Verfahren.

\Needspace{8\baselineskip}
7. Wie zuvor sei

\[
(1+n\cos\varphi)^{\nu}=A+A^{(1)}\cos\varphi+A^{(2)}\cos2\varphi+.\,.\,.\,.+A^{(h)}\cos h\varphi+.\,.\,.
\]

\noindent Bildet man auf beiden Seiten die Differentiale der Logarithmen und kürzt $d\varphi$ heraus, so ergibt sich

\[
\frac{\nu n\sin\varphi}{1+n\cos\varphi}=\frac{A^{(1)}\sin\varphi+2A^{(2)}\sin2\varphi+3A^{(3)}\sin3\varphi+.\,.\,.+hA^{(h)}\sin h\varphi+.\,.\,.}{A+A^{(1)}\cos\varphi+A^{(2)}\cos2\varphi+.\,.\,.+A^{(h)}\cos h\varphi+.\,.\,.}
\]

$\alpha$) Beseitigt man auf beiden Seiten die Nenner durch Multiplikation, so erhält man auf der einen Seite

\[
\begin{aligned}
&\nu n\sin\varphi\bigl(A+A^{(1)}\cos\varphi+A^{(2)}\cos2\varphi+.\,.\,.\,.+A^{(h)}\cos h\varphi\\&\quad+A^{(h+1)}\cos(h+1)\varphi+A^{(h+2)}\cos(h+2)\varphi+.\,.\,.\,.\bigr).
\end{aligned}
\]
